% ============================================================
% preambulo-estetica.tex
% ------------------------------------------------------------
% CAPA REEMPLAZABLE. ES LA QUE UNA EDITORIAL CAMBIA POR LA SUYA
% CUANDO QUIERE MANTENER SU IDENTIDAD VISUAL. NO DEBE DEFINIR
% NINGUNO DE LOS ENTORNOS DEL CONTRATO.
% ------------------------------------------------------------
% ESTE ARCHIVO ES EL ESQUELETO. LOS BLOQUES DE DISEÑO QUE YA
% EXISTEN EN EL preambulo.tex EN USO —titlesec, titletoc,
% estilos de página, diseño de notas al pie, newfloat,
% imakeidx, glossaries, froufrou, siunitx— SE PEGAN EN EL
% APARTADO 8. NO ESTÁN ACÁ PORQUE SON DECISIONES DE DISEÑO
% PROPIAS Y NO CORRESPONDE REESCRIBIRLAS.
% ============================================================
% LAS CUATRO PRESTACIONES
% ------------------------------------------------------------
% El derivado emite PRESTACIONES y no un «modo». Así, agregar
% un estado editorial nuevo no obliga a tocar esta capa, y cada
% bandera dice qué hace en vez de cómo se llama.
%
%   estado_libro     showframe  homeoarchy  enlaces  cruces
%   borrador             si         si         si      no
%   en_correccion        no         si         si      no
%   en_produccion        no         no         si      no
%   publicado            no         no         no      si
%
% LAS CUATRO TOCAN OPCIONES DE PAQUETE, QUE NO SE PUEDEN
% SUPERPONER DESPUES DE CARGADO EL PAQUETE. POR ESO SON
% BANDERAS DEL DERIVADO Y NO REEMPLAZOS DE TEXTO.
% LA EXCEPCION ES hyperref: SE CARGA UNA SOLA VEZ EN EL
% CONTRATO Y LOS COLORES VAN POR \hypersetup, QUE SI SE
% SUPERPONE.
% ============================================================
% CORRECCIONES APLICADAS RESPECTO DEL PREAMBULO EN USO
%   - FUERA \usepackage[utf8]{luainputenc}: LuaLaTeX CONSUME
%     UTF-8 NATIVAMENTE (SC-02).
%   - FUERA \usepackage{svg}: EXIGE Inkscape Y --shell-escape,
%     Y EN LIBROS SOLO SE USAN PNG, JPG Y PDF.
%   - \renewcommand{\footnoterule} ESTABA DEFINIDO DOS VECES
%     CON EL MISMO CUERPO. QUEDA UNO.
%   - babel LO CARGA EL CONTRATO, ANTES DE csquotes Y biblatex,
%     QUE ES EL ORDEN QUE PIDE LA DOCUMENTACION.
%   - luatexja SOLO SI EL LIBRO LO NECESITA.
% ============================================================

\usepackage{ifthen}
\usepackage{calc}

% ============================================================
% 1. GEOMETRIA
% ------------------------------------------------------------
% Los valores vienen del derivado.
%
% showframe ES OPCION DE PAQUETE Y NO SE PUEDE SUPERPONER: SE
% PASA CON \PassOptionsToPackage PARA NO DUPLICAR LA LISTA
% ENTERA DE OPCIONES EN DOS RAMAS DE UN CONDICIONAL.
%
% Hoy los cuatro margenes son iguales en los 45 formatos de la
% base, asi que inner/outer da lo mismo que centering; se usan
% igual para que el dia que se carguen asimetricos no haya que
% tocar nada. En un libro encuadernado el margen interior suele
% necesitar mas aire que el exterior.
% ============================================================
\ifgbshowframe
  \PassOptionsToPackage{showframe}{geometry}
\fi

\usepackage[
  paperwidth=\gbanchopagina,
  paperheight=\gbaltopagina,
  inner=\gbmargeninterior,
  outer=\gbmargenexterior,
  top=\gbmargensuperior,
  bottom=\gbmargeninferior,
  includehead,
  includefoot,
  headsep=14pt,
  footskip=24pt
]{geometry}

% ------------------------------------------------------------
% CRUCES DE CORTE
% ------------------------------------------------------------
% La hoja de montaje es mayor que el formato final: el derivado
% la calcula sumando el margen de cruces por lado. Las cruces
% solo necesitan lugar, y ese lugar es el mismo para cualquier
% formato.
%
% crop DECLARA LA HOJA DE MONTAJE Y geometry EL FORMATO FINAL:
% SON DOS PARES DE MEDIDAS INDEPENDIENTES, NO UNA DERIVADA DE
% LA OTRA. VERIFICADO: crop EXPANDE MACROS EN SUS OPCIONES
% keyval, ASI QUE NO HACE FALTA \crop[...] APARTE.
% ------------------------------------------------------------
\ifgbcruces
  \usepackage[width=\gbanchomontaje,height=\gbaltomontaje,cam,center]{crop}
\fi

% ============================================================
% 2. IDIOMA
% ------------------------------------------------------------
% babel lo carga el contrato, antes de csquotes y biblatex.
% Aca solo se cambian las cadenas, que si son diseno.
% ============================================================
\addto\captionsspanish{\renewcommand{\contentsname}{Sumario}}

% ============================================================
% 3. TIPOGRAFIA
% ------------------------------------------------------------
% Las fuentes son de la estetica, no de la base. Por eso cuando
% una editorial reemplaza esta capa, cambia las fuentes y el
% colofon las nombra correctamente sin que haya que tocar nada
% mas: el contrato solo las lee.
% ============================================================
\usepackage{fontspec}
\usepackage[final,babel=true]{microtype}

\def\gbfuentetexto{Libertinus Serif}
\def\gbfuentetitulos{IBM Plex Sans Condensed}
\def\gbfuentemono{IBM Plex Mono}

\setmainfont{\gbfuentetexto}[
  Numbers={OldStyle,Proportional},
  Ligatures={TeX,Common,Discretionary},Scale=1.18]
\setsansfont{\gbfuentetitulos}[Scale=MatchLowercase,Ligatures=TeX]
\setmonofont{\gbfuentemono}[Scale=0.91]

\usepackage{unicode-math}
\setmathfont{Libertinus Math}[Scale=MatchLowercase]

% SOPORTE CJK SOLO SI EL LIBRO LO NECESITA: luatexja PESA EN
% CADA COMPILACION Y CASI NINGUN LIBRO LO USA.
\ifgbcjk
  \usepackage{luatexja-fontspec}
  \setmainjfont{FandolSong}
\fi

\renewcommand{\normalsize}{\fontsize{10pt}{14pt}\selectfont}
% \topskip=14pt
\normalsize

% ============================================================
% 4. LISTAS, NOTAS Y FILETES
% ============================================================
\usepackage{enumitem}
\setlist{nosep,topsep=4pt}

\usepackage[bottom,stable,hang]{footmisc}
\setlength{\footnotesep}{10pt}
\renewcommand{\footnoterule}{%
  \kern-3pt\hrule height 0.5pt width 0.4\columnwidth\kern6pt}
\renewcommand*{\thefootnote}{\scriptsize\sffamily[\arabic{footnote}]}

% PENDIENTE DE VERIFICAR: EL PREAMBULO EN USO CARGA TAMBIEN
% scrextend Y USA \deffootnote. footmisc CON hang Y \deffootnote
% GOBIERNAN LA MISMA SANGRIA POR MECANISMOS DISTINTOS. HAY QUE
% DECIDIR CUAL QUEDA Y SACAR EL OTRO.

\usepackage[labelfont=bf,font=small,labelsep=period,format=plain]{caption}
\usepackage{needspace}

\setcounter{tocdepth}{1}
\setcounter{secnumdepth}{3}
\renewcommand{\thepart}{\arabic{part}}

\raggedbottom
\clubpenalty=10000
\widowpenalty=10000

% ============================================================
% 5. MARCAS TIPOGRAFICAS DE TRABAJO
% ------------------------------------------------------------
% impnattypo marca homeoarquias y problemas de principio y fin
% de linea. La opcion draft es opcion de paquete y no se puede
% superponer: por eso la carga va condicionada y no comentada
% a mano.
% ============================================================
\ifgbhomeoarchy
  \IfFileExists{impnattypo.sty}{%
    \usepackage[hyphenation,homeoarchy,homeoarchywordcolor=orange,
                homeoarchycharcolor=orange,draft]{impnattypo}}{%
    \PackageWarning{gbpublisher}{impnattypo no instalado:
      se pierden las marcas tipograficas de trabajo}}
  \overfullrule=5pt
\fi

% ============================================================
% 6. ENLACES
% ------------------------------------------------------------
% Se redefine \gbopcioneshyper, que el contrato aplica con
% \hypersetup despues de cargar hyperref. Asi hyperref se carga
% una sola vez y no hace falta duplicar el \usepackage.
%
% En la salida para imprenta los enlaces no deben imprimirse
% coloreados: hidelinks los deja del color del texto y sin
% recuadro, conservando la navegacion en el PDF.
% ============================================================
\usepackage{xcolor}
\definecolor{gbenlace}{RGB}{160,0,120}

\ifgbenlaces
  \renewcommand{\gbopcioneshyper}{colorlinks=true,
    linkcolor=gbenlace,citecolor=gbenlace,urlcolor=gbenlace,
    filecolor=gbenlace,linktoc=all}
\else
  \renewcommand{\gbopcioneshyper}{hidelinks,pdfborder={0 0 0}}
\fi

% ============================================================
% 7. CONSTANCIA EN EL LOG
% ------------------------------------------------------------
% Deja escrito en el .log con que prestaciones se compuso.
% Cuando un PDF sale distinto de lo esperado, esta linea dice
% por que sin tener que abrir el derivado.
% ============================================================
\typeout{[gbpublisher] prestaciones:
  showframe=\ifgbshowframe si\else no\fi,
  homeoarchy=\ifgbhomeoarchy si\else no\fi,
  enlaces=\ifgbenlaces si\else no\fi,
  cruces=\ifgbcruces si\else no\fi}

% ============================================================
% 8. AQUI VAN LOS BLOQUES DE DISENO DEL PREAMBULO EN USO
% ------------------------------------------------------------
%   - titlesec y los estilos de titulo
%   - titletoc: partes, capitulos y secciones del sumario
%   - titleps: estilos de pagina y folios corridos
%     (\gbtituloabreviado ES LO QUE CORRESPONDE EN EL FOLIO
%      CORRIDO, Y HOY EL PREAMBULO NO LO CONSUME)
%   - newfloat: el entorno flotante "imagen"
%   - imakeidx + xindy: names (lo llena esindex desde biblatex),
%     onomastico y concepto (los carga el editor)
%   - glossaries: siglas y glosario
%   - siunitx, froufrou, qrcode, pdfpages, rotating
% ============================================================
